\subsection{Exact extraction of the centered $m$ event}
\label{component:p8da9ea6ed9e5}
\label{p8da9ea6ed9e5:sec:r17-r5}

The identities in this subsection concern the entire four-state parent
\eqref{eq:c2-ba-parent}. They do not transfer a source admission from one
stratum to another. All source quantities belong to $K_c$, all testing
quantities belong to $K_t$, and the two labelled factors remain separate.
Put $H=hB$, $a=x/B$ and $g=1-a(1-r-H)$, as in the original stationary
source. Write
\[
 M=vW-X+uvY-uZ,\qquad
 \xi=X+ruZ,\quad\eta=Y-huW,\quad
 R_2=vW-a^{-1}uZ,\quad T=uvW-BZ.
\]
Direct expansion gives the contact-centered row
\begin{equation}
 R_{mc}=M-gvW=-\xi+uv\eta+(1-g)R_2+huT.
 \tag{R5.2}\label{eq:c2-r55-centered-row}
\end{equation}
Indeed $1-g=a(1-r-H)$ and $H=hB$. At the original contact the four
roots on the right vanish. Thus retaining a multiple of $R_{mc}^2$
preserves the contact value and its full first jet.

For any positive real $K<m$, subtract this whole square from the
original parent. With $P$ denoting \eqref{eq:c2-ba-parent}, the exact
remainder is
\[
 \begin{split}
 P-(m-K)R_{mc}^2
   &=P\big|^{\rm event}_{m=K}
       +2\tau vWM-\tau g v^2W^2,\\
 \tau&=(m-K)g.
 \end{split}
\]
Here $P|^{\rm event}_{m=K}$ means changing only the coefficient of
$M^2$ in the displayed event parent; all other actual source
coefficients remain fixed. Equivalently, its complete $m$ contribution
is the elementary identity
\begin{equation}
 mM^2-(m-K)R_{mc}^2
   =KM^2+2(m-K)g\,vWM-(m-K)g^2v^2W^2.
 \tag{R5.5}\label{eq:c2-r55-event}
\end{equation}
This is an identity for arbitrary $W,X,Y,Z$, before restricting a
state or taking a residue. In particular, when $g$ is infinitesimal,
$mg$ is bounded and the other event coefficients are bounded, its
remainder has bounded coefficients at every bounded testing pair.
It does not require $a$ to be bounded. The retained term is a square
in the tensor ring because $\sqrt{m-K}$ belongs to the source field.

The same extraction has an exact five-root matrix form. In the
unnormalized WA coordinates
${\cal R}=(\xi,\eta,R_1,R_2,T)^T$, where $R_1=uW-t^{-1}vZ$,
let $G$ be the full WA matrix, including both complete normal rows
and its hard block. Put $L_m=(-1,uv,0,1-g,hu)^T$. Then
\begin{equation}
 G=G\big|^{\rm table}_{m=K}+(m-K)L_mL_m^T+\tau H_m,
 \tag{R5.3}\label{eq:c2-r55-matrix}
\end{equation}
where the only nonzero entries of the symmetric matrix $H_m$ are
\begin{equation}
 (H_m)_{14}=-1,\qquad (H_m)_{25}=v,\qquad
 (H_m)_{44}=2-g,\qquad (H_m)_{45}=hu.
 \tag{R5.4}\label{eq:c2-r55-correction}
\end{equation}
The full WA table is affine in $m$. Its entrywise derivative, with
every other actual source coefficient fixed, is
$\partial_mG=L_mL_m^T+gH_m$, which proves
\eqref{eq:c2-r55-matrix}. This formal table substitution differs
from the event substitution in \eqref{eq:c2-r55-event}: the WA table
has already used the original contact balances. Neither substitution
asserts that a modified source is physical or stationary. Both
expressions describe the same original parent after removing
$(m-K)R_{mc}^2$.

\paragraph{A real cap on a squared row.}
\label{lem:c2-r55-real-cap}
Let $B_0$ be a real symmetric form and let $M_0$ be a nonzero real
linear row. If $B_0$ is positive definite on $\ker M_0$, then
$B_0+KM_0^2$ is positive definite for all sufficiently large positive
real $K$. To see this, choose real coordinates $(z,q)$ with
$q=M_0$ and write the form as
\begin{equation}
 z^TAz+2q\,b^Tz+(c+K)q^2,\qquad A\succ0.
 \label{eq:c2-r55-real-cap}
\end{equation}
Its real Schur complement is $c+K-b^TA^{-1}b$, which is positive
once $K$ is sufficiently large. This uses only real matrices. The
argument also applies after adding a fixed real multiple of $LM_0$
for any real row $L$, because that addition vanishes on $\ker M_0$.

\paragraph{The actual unbounded-$a$ R6 finite remote source.}
Retain the original physical R6 source, its stationary contact, all
strict gaps, coefficient signs, total scale and prior exclusions.
Thus $t,h,l=1-h,j,r$ are positive bounded units,
$B,x=aB,y=tB,g$ are positive infinitesimals, $a$ is unbounded,
and $n\ge0$. No lower bound on $b=a^2B$ is imposed. Put
$s=a^{-1}$ and $w=a(1-r)=1+hx-g$, so $w\sim1$,
$x=sb$ and $B=s^2b$. The physical identity
\eqref{eq:c2-ba-physical} gives
\[
 2r\Delta_*n/s^2=wQ_*-bK_*,\qquad
 \Delta_*=s\Delta\sim1,\quad Q_*\sim1,\quad
 K_*=(1+s)(h+t)(1+x)\sim h+t.
\]
Here $Q_*$ is exactly $Q_s$ in that identity. Since $n\ge0$, this
forces $b$ to be bounded. Substitution in the actual rational source
formulas \eqref{eq:c2-ba-source}, followed by cancellation before
taking residues, yields jointly
\begin{equation}
 \begin{gathered}
 (mg)_0=\tfrac12,\qquad n_0=\alpha_0=0,\\
 A_0={t_0\over2h_0},\qquad
 C_0={h_0+t_0\over2l_0},\qquad
 D_0={1\over2(t_0-l_0)},\\
 (n/s^2)_0={1-b_0(h_0+t_0)\over2},\qquad
 \left({\alpha\over s^3b^2}\right)_0={h_0+t_0\over2}.
 \label{eq:c2-r55-r6-source}
 \end{gathered}
\end{equation}
The canceled denominators are units jointly at $s=g=0$ for bounded
$b$. In particular these formulas cover $b$ a unit or infinitesimal,
and literal $n=0$. The positive unit $j$ gives $t_0>l_0$, while
$h_0,l_0>0$. Thus $A_0,C_0,D_0$ are strictly positive and $m$ is
unbounded. No estimate compares $s$, $g$ or $b$ with a testing variable.

Consider a bounded testing pair with $v_0\ne0$, and write $u,v,h,t,l$
for its real testing residues and the real source residues in this
paragraph. Set $U=u^2$, $V=v^2$, $D_u=1+U$, $D_v=1+V$.
By \eqref{eq:c2-r55-event}, the complete bounded remainder has residue
\[
 B_0+KM_0^2+vWM_0,\qquad M_0=vW-X+uvY-uZ,
\]
where the full four-state form $B_0$ is
\[
 \begin{split}
 B_0={}&C_0D_v(uW-Y)^2+A_0D_uD_vY^2
       +(D_0+\tfrac12)(uW+uvX-Y-vZ)^2\\
      &+\tfrac12D_u(vY-Z)^2
       +\tfrac12D_v(X+uZ)^2-\tfrac12D_uD_vZ^2.
 \end{split}
\]
On $M_0=0$, substitute $X=vW+uvY-uZ$. The determinant of this
three-state form in $(W,Y,Z)$ is exactly
\begin{equation}
 {D_u^2D_v^2(h+t)\over8hl(t-l)}
       \{V^2+t(h+t)U^2V+hUV^2\}>0.
 \label{eq:c2-r55-r6-kernel}
\end{equation}
At $Z=0$ its signed core is
$[J^2+D_uVY^2+D_vX^2]/2$, with
$J=uW+uvX-Y$. The positive $A_0$ term forces $Y=0$,
the $X^2$ term then forces $X=0$, and $M_0=0$ with $v\ne0$
forces $W=0$. Hence the restriction to $Z=0$ is positive definite.
Together with the determinant this proves $B_0|_{M_0=0}\succ0$.
The centered correction $vWM_0$ vanishes on that same kernel.
The real-cap argument above therefore gives a positive real $K$
for which the full four-state residue is strictly positive definite.

Unbounded actual $m$ permits this real choice with $K<m$. All
coefficients of the exact remainder belong to
$B_c\otimes_{\mathbb R}B_t$: the source coefficients
$n,A,C,D,\alpha,g,\tau$ are bounded, and $u,v$ are bounded in their
own field. Lemma~\ref{p25af2c58da7e:lem:gram} now gives its exact
finite tensor sum of squares for arbitrary state coordinates.
Restoring the retained square proves precisely this finite remote
join. The same real identity applies separately to real and imaginary
state parts. The projective endpoint identities follow in the next subsection.

\subsection{Exact projective full-state identities for R6 and R7}
\label{component:pd5fa361b4623}
\label{pd5fa361b4623:sec:r17-r6-bounded-a}

The following polynomial identities concern the original parent
\eqref{eq:c2-master-parent}, rather than a source theorem on another
stratum. They also give an endpoint proof that applies when the actual
coefficient $n$ is zero or the actual coefficients $A,D$ are unbounded.
Write $D_z=1+z^2$ for each angular variable $z$, and continue to use
$P=F/4$. Thus
\[
\begin{split}
 P={}&nD_uD_vX^2+mM^2+CD_v(uW-Y)^2+AD_uD_vY^2+DJ^2
       +\alpha D_u(vY-Z)^2\\
 &+\tfrac12\{J^2+D_u(vY-Z)^2+D_v(X+uZ)^2-D_uD_vZ^2\},\\
 M={}&vW-X+uvY-uZ,\qquad J=uW+uvX-Y-vZ.
\end{split}
\]

\paragraph{Three whole homogeneous polynomials.}
For $p=1/u$ define
\[
 M_u=pvW-pX+vY-Z,\qquad J_u=W+vX-pY-pvZ.
\]
Multiplication by $p^2$ gives the exact identity
\begin{equation}\label{eq:c2-r65-pu}
\begin{split}
 P_u:={}&p^2P
 =nD_pD_vX^2+mM_u^2+CD_v(W-pY)^2+AD_pD_vY^2\\
 &+DJ_u^2+\alpha D_p(vY-Z)^2+\tfrac12\mathcal C_u,\\
 \mathcal C_u={}&J_u^2+D_p(vY-Z)^2+D_v(pX+Z)^2-D_pD_vZ^2.
\end{split}
\end{equation}
For $q=1/v$ put
\[
 M_v=W-qX+uY-quZ,\qquad J_v=quW+uX-qY-Z.
\]
The other single homogeneous polynomial is
\begin{equation}\label{eq:c2-r65-pv}
\begin{split}
 P_v:={}&q^2P
 =nD_uD_qX^2+mM_v^2+CD_q(uW-Y)^2+AD_uD_qY^2\\
 &+DJ_v^2+\alpha D_u(Y-qZ)^2+\tfrac12\mathcal C_v,\\
 \mathcal C_v={}&J_v^2+D_u(Y-qZ)^2+D_q(X+uZ)^2-D_uD_qZ^2.
\end{split}
\end{equation}
For both reciprocal coordinates use the five roots
$(X,W,Y,P_Z,Q_Z)$, where $P_Z=pZ$ and $Q_Z=qZ$, and set
\[
 M_b=pW-pqX+Y-Q_Z,\qquad J_b=qW+X-pqY-P_Z.
\]
Then the complete polynomial is
\begin{equation}\label{eq:c2-r65-pb}
\begin{split}
 P_b:={}&p^2q^2P
 =nD_pD_qX^2+mM_b^2+CD_q(W-pY)^2+AD_pD_qY^2\\
 &+DJ_b^2+\alpha D_p(Y-Q_Z)^2+\tfrac12\mathcal C_b,\\
 \mathcal C_b={}&J_b^2+D_p(Y-Q_Z)^2
                  +D_q(p^2X^2+2XP_Z-P_Z^2).
\end{split}
\end{equation}
In particular the apparently remaining $Z^2$ term has canceled in
\eqref{eq:c2-r65-pb}. This is an identity on the five displayed roots;
their independence is unnecessary. All three right-hand sides are
polynomials at $p=0$ or $q=0$. They therefore define the literal
projective endpoints without division by a reciprocal coordinate.

\paragraph{The R6 kernel identities, including zero residue of $n$.}
Suppose $g$ is infinitesimal, $m$ is unbounded, $mg$ is bounded,
$n,C,A,D,\alpha$ are bounded, $n\ge0$, and
$C_0,A_0,D_0>0$, $\alpha_0=0$. These are precisely the coefficient
properties needed here. Both the original bounded-$a$ R6 source
and the unbounded-$a$ R6 sources supply them; in the latter,
$n_0=0$, including when $n$ is literally zero.

For any positive real $K<m$, apply the exact identity
\begin{equation}\label{eq:c2-r65-centered}
 mH^2-(m-K)(H-gL)^2
   =KH^2+2\mu LH-\mu gL^2,\qquad \mu=(m-K)g,
\end{equation}
with $(H,L)=(M_u,pvW),(M_v,W),(M_b,pW)$ in the three
homogeneous polynomials. The remainder has bounded coefficients.
Write $\kappa=\mu_0=(mg)_0$.

At $p=0$ the centered correction has zero residue. On the kernel
$vY-Z=0$ of the limiting $m$ row, the complete residue is
\begin{equation}
 n_0D_vX^2+C_0D_vW^2+A_0D_vY^2
       +(D_0+\tfrac12)(W+vX)^2.\tag{R6.5}\label{eq:c2-r65-kernel-u}
\end{equation}
Here and in the next two displays the remaining testing letters denote
their real residues. The full four-state residue is this form plus
$(K+\tfrac12)(vY-Z)^2$. If $n_0>0$, it is strict for every finite
$v$. If $n_0=0$, the determinant of the kernel form is
\[
 A_0C_0(D_0+\tfrac12)v^2D_v^2.
\]
It is strict for $v\ne0$. At $v=0$ its $X$ direction is absent;
the separate four-root calculation below treats that corner.

At $q=0$ the centered correction is
$2\kappa W(W+uY)$, which vanishes on the limiting $m$-row kernel
$W+uY=0$. On this kernel the complete residue is
\begin{equation}
 (n_0+\tfrac12)D_uX^2+C_0D_u^2Y^2
       +(A_0+\tfrac12)D_uY^2+D_0(uX-Z)^2.\tag{R6.6}\label{eq:c2-r65-kernel-v}
\end{equation}
Its determinant on $(X,Y,Z)$ is
\[
 D_0(n_0+\tfrac12)D_u^2(C_0D_u+A_0+\tfrac12)>0.
\]
Its $Y$ coefficient is positive and its $(X,Z)$ block is strict, so
the kernel form is positive definite for every finite $u$, including
$u=0$ and $n_0=0$. The real-cap argument in
Section~\ref{lem:c2-r55-real-cap} chooses a sufficiently large real
$K$ for the full four-state residue; its inverses occur only in a
real matrix, never in the tensor ring.

At $p=q=0$ the centered correction again has zero residue and the
complete five-root form is
\begin{equation}
 (n_0+\tfrac12)X^2+C_0W^2+A_0Y^2+D_0(X-P_Z)^2
                 +(K+\tfrac12)(Y-Q_Z)^2.\tag{R6.7}\label{eq:c2-r65-double}
\end{equation}
The two nontrivial block determinants are
$D_0(n_0+\tfrac12)$ and $A_0(K+\tfrac12)$. Both remain positive
at $n_0=0$; the whole five-root determinant is
$A_0C_0D_0(n_0+\tfrac12)(K+\tfrac12)>0$.

\paragraph{A common capped proof for the R6 and R7 endpoints.}
There is also an endpoint proof using only positive real lower bounds
on $m,C,A,D$. It avoids imposing boundedness on $A$ or $D$.
Choose positive real numbers
\[
 m_c<m,\qquad C_c<C,\qquad A_c<A,\qquad D_c<D.
\]
Retain the whole $n,\alpha$ event squares and the four nonnegative
excess event squares with coefficients
$m-m_c,C-C_c,A-A_c,D-D_c$. In each homogeneous chart the remaining
polynomial is exactly \eqref{eq:c2-r65-pu},
\eqref{eq:c2-r65-pv}, or \eqref{eq:c2-r65-pb} with
\[
 (n,\alpha,m,C,A,D)=(0,0,m_c,C_c,A_c,D_c).
\]
Thus all its coefficients are bounded pure-testing elements. Its
complete residues in the three charts are, respectively,
\begin{align}
 &C_cD_vW^2+A_cD_vY^2+(D_c+\tfrac12)(W+vX)^2
                   +(m_c+\tfrac12)(vY-Z)^2,
                         \label{eq:c2-r65-cap-u}\\
 &m_c(W+uY)^2+C_c(uW-Y)^2+(A_c+\tfrac12)D_uY^2
                   +D_c(uX-Z)^2+\tfrac12D_uX^2,
                         \label{eq:c2-r65-cap-v}\\
 &C_cW^2+A_cY^2+D_c(X-P_Z)^2+\tfrac12X^2
                   +(m_c+\tfrac12)(Y-Q_Z)^2.
                         \label{eq:c2-r65-cap-b}
\end{align}
The square rows in \eqref{eq:c2-r65-cap-u} have common kernel zero
when $v\ne0$. The square rows in \eqref{eq:c2-r65-cap-v} have
common kernel zero for every finite $u$: first $X=Y=0$, then the
$m_c$ and $D_c$ rows give $W=Z=0$. In
\eqref{eq:c2-r65-cap-b}, the separate $W,X,Y$ squares and the last
two rows force all five roots to vanish. These are strict Gram
matrices for the entire homogeneous state polynomial.

The R6 bounds give the required positive real lower bounds. In
unbounded-$a$ R7, $m,A$ are unbounded and $C_0=t_0/2>0$.
When $t_0>1$, $D_0=1/[2(t_0-1)]>0$; when $t_0=1$, $D$ is
unbounded. Hence these same real caps exist on both R7 gap
boundaries. The actual $n,\alpha$ packets were retained, so neither
of them was replaced by its zero residue in the full polynomial.

\paragraph{The complete $(\infty,0)$ allocation.}
Assume $p,v$ are infinitesimal, allowing either to be literally zero.
Retain all the original event squares except positive real caps
$C_c,A_c$ of the $C,A$ events. Put
\[
 H=pX+Z,\qquad T=vX,\qquad
 \widehat J=W+D_pT-pY-pvH.
\]
The entire selected polynomial has the following four-root identity:
\begin{equation}\label{eq:c2-r65-exceptional}
\begin{split}
 Q_{\infty0}={}&C_cD_v(W-pY)^2+A_cD_pD_vY^2
                      +\tfrac12\widehat{\mathcal C},\\
 \widehat{\mathcal C}={}&\widehat J^2+D_pv^2Y^2-2D_pvYH
          +(1-p^2v^2)H^2\\
 &+2pD_pTY+2pvD_pTH-p^2D_pT^2.
\end{split}
\end{equation}
Indeed substitute $Z=H-pX$ into $\mathcal C_u$ and collect every
remaining $vX$ as $T$. This gives
\eqref{eq:c2-r65-exceptional} by multiplication and addition only.
In particular it gives every state cross term, with no inverse of
$p$ or $v$. On $(H,W,Y,T)$ its residue is
\[
 \tfrac12H^2+\tfrac12(W+T)^2+C_cW^2+A_cY^2\succ0,
 \qquad \det=\tfrac14A_cC_c>0.
\]
This replaces the nonstrict ordinary state residue at $n_0=v_0=0$.
It also treats the literal endpoint, where some of the four roots
are dependent or zero.

\paragraph{Passage from the residues to tensor sums of squares.}
Each polynomial above has an explicit finite root frame with bounded
coefficients. For the centered R6 formulas the coefficients belong to
$B_c\otimes_{\mathbb R}B_t$ because $\mu=(m-K)g$ is bounded;
for the capped proof they even belong to $B_t$. The strict real
matrices therefore meet every hypothesis of
Lemma~\ref{p25af2c58da7e:lem:gram}. That lemma gives a finite tensor
sum of squares for the complete remainder, rather than merely
pointwise real positivity. Restore the retained positive event
squares, or the positive centered excess square in
\eqref{eq:c2-r65-centered}. For nonzero reciprocal coordinates,
restoring homogeneity multiplies squares only by pure-testing
factors. At literal endpoints the displayed polynomial value is
used directly. All identities hold for either angle sign and every
real state; a complex state is treated by adding the identities for
its real and imaginary parts. These calculations cover every
projective test with at least one angle at infinity or unbounded,
without using a finite-contact theorem on another source stratum.
